\subsubsection{Hardy--inner factorization horizon semantics: leakage--offset sum, Poisson phase law, and Clark/Carleson certificates}\label{subsubsec:xi-hardy-inner-clark-carleson}
This paragraph further compresses the semantic conditions of "horizon zero-knowledge--unitary slice" into inner-outer structural constraints in the Hardy/Smirnov factorization, and provides three classes of auditable consequences:
\textbf{(i)} critical line pure-phase property forces the completed determinant to be an inner function in the right half-plane, so that the outer factor (amplitude leakage) vanishes;
\textbf{(ii)} the amplitude leakage of approximate unitarity can quantitatively upper-bound the weighted total offset of off-line zeros, and write the phase derivative as a Poisson (Breit--Wigner) kernel superposition;
\textbf{(iii)} with Clark measures and Carleson/BMO as intermediaries, elevate "worst-case soundness/recovery rate parametrized by challenge seed" to a geometric measure criterion of zero density.
\par
This paragraph does not introduce new empirical assumptions; it merely repurposes the factorization language of Appendix \S\ref{subsec:unit-disk-inner-outer} and Subsubsection \ref{subsubsec:selfdual-blaschke-defect-fixed-slice} under the bounded type (or Smirnov) regime, and switches between half-plane and disk formulations via the Cayley transform.

\paragraph{Half-plane regime and disk notation}
Denote the right half-plane by
\[
\mathbb{H}_{1/2}:=\{s\in\mathbb{C}:\ \Re(s)>\tfrac12\}.
\]
Suppose \(D(s)\) is analytic on \(\mathbb{H}_{1/2}\) and of bounded type, and satisfies the reflection symmetry
\begin{equation}\label{eq:xi-inner-reflection-symmetry}
\boxed{\;D(s)=\overline{D(1-\overline{s})}\;}
\end{equation}
(this form is isomorphic to the determinant reflection symmetry of Proposition \ref{prop:xi-hzom-functional-equation}).
Let the shifted variable \(w:=s-\tfrac12\) send \(\mathbb{H}_{1/2}\) to the standard right half-plane \(\Re(w)>0\);
then apply the Cayley transform to send it to the unit disk \(\mathbb{D}\), obtaining the disk-ified function \(\mathfrak{D}\) (defined up to composition isomorphism), and take its nontangential boundary value representative on \(\TT=\partial\mathbb{D}\).

\begin{theorem}[Hardy--inner characterization: horizon zero-knowledge unitarity eliminates outer factor]\label{thm:xi-hardy-inner-characterization}
Under the above setting, if the "no amplitude leakage" condition of horizon zero-knowledge holds on the critical line, i.e., for almost every \(t\in\mathbb{R}\) we have
\[
\boxed{\ |D(\tfrac12+it)|=1\ },
\]
then \(D\) is an inner function on \(\mathbb{H}_{1/2}\): its Hardy regular boundary value exists and belongs to \(L^\infty(\mathbb{R})\).
Equivalently, the disk-ification \(\mathfrak{D}\) is a disk inner function and has the unique (canonical) factorization
\[
\boxed{\;
\mathfrak{D}(z)=e^{i\beta}\,B(z)\,S(z)\;},
\]
where \(B\) is the Blaschke product formed by zeros inside the disk, \(S\) is the singular inner factor, and the outer factor degenerates to a constant phase factor of modulus \(1\).
\end{theorem}
\begin{proof}[Proof outline]
A disk-ified function of bounded type has the canonical inner-outer decomposition \(\mathfrak{D}=B\,S\,O\) (Appendix \S\ref{subsec:unit-disk-inner-outer}).
Critical line pure-phase property is equivalent to \(|\mathfrak{D}(e^{i\theta})|=1\) almost everywhere, so the outer factor \(O\) determined by the boundary modulus must be a constant phase.
Thus \(\mathfrak{D}\) is an inner function and can be written as \(e^{i\beta}BS\). This completes the proof.
\end{proof}

\begin{proposition}[Leakage--zero offset sum: outer factor logarithmic modulus upper-bounds off-line zero offset]\label{prop:xi-leakage-zero-offset-sum}
In the generalized case of Theorem \ref{thm:xi-hardy-inner-characterization}, do not assume \(|D(\tfrac12+it)|=1\).
Let
\[
L:=\frac{1}{2\pi}\int_{-\infty}^{\infty}\log |D(\tfrac12+it)|\,\frac{\dd t}{1+t^2},
\]
and assume \(\log|D(\tfrac12+it)|\in L^1((1+t^2)^{-1}\dd t)\).
Then the set of off-line zeros in the right half-plane \(\{\rho\in\mathbb{H}_{1/2}:D(\rho)=0\}\) satisfies the offset sum inequality
\[
\boxed{\;
\sum_{\rho:\Re(\rho)>\frac12}\frac{\Re(\rho)-\frac12}{1+(\Im(\rho))^2}\ \le\ L.\;}
\]
Moreover, equality holds if and only if the singular inner factor is trivial and the outer factor is completely determined by the critical line modulus.
\end{proposition}
\begin{proof}[Proof outline]
Disk-ify \(D\) to \(\mathfrak{D}\) and take its inner-outer decomposition \(\mathfrak{D}=BSO\).
The Poisson integral of the logarithmic modulus at the base point gives the harmonic value of \(\log|O|\); and the Blaschke condition relates the sum of distances from zeros to the boundary to this harmonic value.
Choosing the base point corresponding to the weight \((1+t^2)^{-1}\) on \(\mathbb{H}_{1/2}\) yields the stated inequality. This completes the proof.
\end{proof}

\begin{proposition}[Poisson superposition law of critical line phase derivative: zero contributions are Breit--Wigner kernels]\label{prop:xi-phase-derivative-poisson-superposition}
Still in the bounded type regime, let \(\Theta(t):=\arg D(\tfrac12+it)\) be taken as a differentiable representative (in the almost-everywhere sense).
Then there exists a non-negative finite measure \(\mu\) (corresponding to the singular inner factor) such that for almost every \(t\) we have
\[
\boxed{\;
\Theta'(t)=
\sum_{\rho:\Re(\rho)>\frac12}
\frac{\Re(\rho)-\frac12}{(t-\Im(\rho))^2+(\Re(\rho)-\frac12)^2}
\ +\ \frac{\dd\mu}{\dd t}(t)\;}
\]
where the sum converges by the Blaschke condition, and \(\dd\mu/\dd t\) denotes the absolutely continuous part density of \(\mu\) (if it exists).
Therefore each off-line zero produces a standard Poisson (Breit--Wigner) kernel profile on the critical line; if the singular inner factor is trivial, then \(\Theta'\) is completely given by the Poisson superposition of off-line zeros.
\end{proposition}
\begin{proof}[Proof outline]
For the disk-ified inner function \(\mathfrak{D}=BS\), differentiate the boundary phase: the Blaschke factor gives the Poisson kernel superposition, and the singular factor gives the boundary derivative term of the singular measure.
Change variables back to the \(t\)-coordinate to obtain the conclusion. This completes the proof.
\end{proof}

\begin{corollary}[Horizon semantics of Clark measure: spectral measure parametrized by challenge phase seed]\label{cor:xi-clark-measure-horizon-semantics}
Let \(\mathfrak{D}\) be the disk-ified inner function of Theorem \ref{thm:xi-hardy-inner-characterization}.
Then for each \(\alpha\in\TT\) there exists a unique Aleksandrov--Clark measure \(\sigma_\alpha\) such that for all \(z\in\mathbb{D}\) we have
\[
\boxed{\;
\Re\!\left(\frac{\alpha+\mathfrak{D}(z)}{\alpha-\mathfrak{D}(z)}\right)
=\int_{\TT} P_z(\zeta)\,\dd\sigma_\alpha(\zeta)\;}
\]
where \(P_z\) is the Poisson kernel.
If we interpret \(\alpha\) as the verifier's challenge "phase seed", then:
\begin{enumerate}
  \item The singular part of \(\sigma_\alpha\) is non-zero if and only if the singular inner factor is non-trivial (there exists a "spectral spike/zero-measure spike leakage" mechanism);
  \item Atoms of \(\sigma_\alpha\) correspond to critical-line phase stationary points (equivalent to stable solutions of \(\mathfrak{D}(e^{i\theta})=\alpha\)), which can be viewed as spectralized readout of worst-case soundness situations.
\end{enumerate}
\end{corollary}

\begin{theorem}[Conservative system synthesis: inner function \(\Longleftrightarrow\) horizon unitary dilation]\label{thm:xi-inner-iff-conservative-colligation}
Suppose \(\mathfrak{D}\) is a disk inner function.
Then there exists a conservative (unitary scattering system in the unitary conjugacy class of a colligation) \(\mathfrak{U}\) such that \(\mathfrak{D}\) is its transfer function; conversely, the transfer function of any conservative system must be an inner function.
Therefore under the abstract encoding of "horizon interface = scattering, transcription = output", perfect zero-knowledge unitarity is strictly equivalent to "existence of a unified unitary dilation realizing all challenge--response dynamics".
\end{theorem}

\begin{proposition}[Zero-knowledge entropy production \(\equiv\) outer factor: reversibility forces outer factor to vanish]\label{prop:xi-entropy-production-equals-outer-factor}
Under the bounded type regime, the disk-ified function has the inner--outer decomposition \(\mathfrak{D}=I\cdot O\).
Its outer factor satisfies the harmonic extension formula: for almost every boundary point,
\[
\boxed{\;
\log|O(e^{i\theta})|
=\mathcal{P}\!\bigl(\log|\mathfrak{D}(e^{i\cdot})|\bigr)(\theta)\;}
\]
where \(\mathcal{P}\) is the Poisson integral operator.
If we interpret \(\log|\mathfrak{D}(e^{i\theta})|\) as "entropy production density relative to the visible simulator", then the outer factor is completely equivalent to the harmonic extension of this entropy production;
in particular, zero entropy production (strict reversibility/detailed balance) holds if and only if the outer factor is constant modulus \(1\), so that \(\mathfrak{D}\) degenerates to a pure inner function and satisfies pure-phase property on the critical line.
\end{proposition}

\begin{proposition}[Off-critical zeros force transcript length lower bound: resonance offset \(\Rightarrow\) minimum information complexity of near-zero-knowledge]\label{prop:xi-offcritical-zero-forces-transcript-length}
Suppose \(D\) has an off-line zero \(\rho=\tfrac12+\delta+i\gamma\) (\(\delta>0\)) in \(\mathbb{H}_{1/2}\).
Assume a family of test statistics on scale parameter \(T\) induced by \(D\) can be realized by transcripts of length at most \(m(T)\),
and require the protocol family to satisfy total variation \(\varepsilon\)-near-zero-knowledge.
Then there exists a constant \(c>0\) such that when \(T\asymp|\gamma|\) we must have
\[
\boxed{\;
m(T)\ \ge\ c\,\delta^{-1}\log\!\frac{1}{\varepsilon}\;}
\]
otherwise one can construct an amplifiable distinguisher from the Poisson resonance term of Proposition \ref{prop:xi-phase-derivative-poisson-superposition}, making the simulated and real distributions incompatible at the \(\varepsilon\) threshold.
\end{proposition}
\begin{proof}[Proof outline]
The Poisson kernel term induced by the off-line zero on the critical line has an amplifiable structure with peak width \(\asymp\delta\) and peak height \(\asymp\delta^{-1}\);
if the transcript length is insufficient to resolve this peak at scale \(T\), any simulator will produce incompressible bias in TV distance.
Aligning this bias with the \(\varepsilon\) constraint yields the \(m(T)\) lower bound of \(\delta^{-1}\log(1/\varepsilon)\) type. This completes the proof.
\end{proof}

\begin{theorem}[Equivalence of zero density and recovery rate: Carleson/BMO criterion template]\label{thm:xi-carleson-bmo-recovery-template}
Suppose \(D\) is of bounded type and satisfies the reflection symmetry \eqref{eq:xi-inner-reflection-symmetry}.
Quantify the near-zero-knowledge recovery rate as: there exists a family of recovery maps such that the error satisfies \(\varepsilon(T)\le e^{-\kappa T}\) (for some \(\kappa>0\)).
Then the following propositions constitute a standard equivalence template (constants depend only on normalization):
\begin{enumerate}
  \item The offset measure of off-line zeros in the right half-plane
  \[
  \nu:=\sum_{\rho:\Re(\rho)>\frac12}\bigl(\Re(\rho)-\tfrac12\bigr)\,\delta_{\rho}
  \]
  is a Carleson measure, with Carleson norm \(\lesssim \kappa\);
  \item The critical line logarithmic modulus \(\log|D(\tfrac12+it)|\) belongs to BMO, with BMO norm \(\lesssim \kappa\);
  \item Zero-free strip width, off-line zero density, and recoverable exponential rate \(\kappa\) mutually control each other on the same scale.
\end{enumerate}
\end{theorem}

\begin{remark}[Protocol interpretation of singular inner factor: zero-measure but amplifiable "information spikes"]\label{rem:xi-singular-inner-factor-information-spikes}
In the inner function decomposition \(\mathfrak{D}=e^{i\beta}BS\) of Theorem \ref{thm:xi-hardy-inner-characterization},
the singular inner factor \(S\) is non-trivial if and only if there exists a family of critical-line zero-measure sets \(E\subset\mathbb{R}\) such that the boundary phase of \(\mathfrak{D}\) exhibits spikes/jumps on \(E\) that cannot be approximated by any \(L^2\)-visible statistical smoothing.
In protocol semantics, this corresponds to "zero-measure but amplifiable" anomalous transcript event families: the simulator cannot uniformly approximate in any operational distance using only horizon data.
Therefore excluding the singular inner factor can be viewed as a strengthened axiom of zero-knowledge at the zero-measure anomalous event level; once this axiom holds, all non-trivial structure is determined by Blaschke zeros, and the zero--phase relationship is completely characterized by the Poisson superposition law of Proposition \ref{prop:xi-phase-derivative-poisson-superposition}.
\end{remark}

\begin{proposition}[Analytic discrimination of continuous hair: real analyticity of boundary density \(\Longleftrightarrow\) no singular inner factor (finite point defect regime)]\label{prop:xi-analytic-boundary-density-iff-no-singular-inner}
Under the setting of Theorem \ref{thm:xi-hardy-inner-characterization}, suppose the disk-ified inner function decomposes as
\(
\mathfrak{D}=e^{i\beta}BS
\).
Let \(\mu\) be the Herglotz measure corresponding to \(S\), and perform the Lebesgue decomposition
\(
\dd\mu=\dd\mu_{\mathrm{ac}}+\dd\mu_{\mathrm{s}}+\sum_j m_j\delta_{\theta_j}
\).
Then:
\begin{enumerate}
  \item \textbf{(Atoms \(\Rightarrow\) Lorentzian peaks)}
  The boundary density induced by the Blaschke/atomic part in the Cayley chart is a finite or countable superposition of Lorentzian (Poisson/Breit--Wigner) kernels,
  thus giving a real analytic function in the "finite point defect" case (see the rational pole--residue inversion of Proposition \ref{prop:xi-phase-derivative-poisson-superposition} and Theorem \ref{thm:xi-offcritical-poisson-herglotz-inversion}).
  \item \textbf{(Singular inner factor \(\Rightarrow\) singular distribution)}
  If \(\mu_{\mathrm{s}}\neq 0\), then the boundary phase density (as a distribution) contains an irreducible singular part, and thus cannot be represented by a real analytic function on any non-empty open interval.
  \item \textbf{(Analyticity discrimination)}
  Therefore, in the finite point defect regime,
  "boundary density is a real analytic function on \(\RR\)" is strictly equivalent to "singular inner factor is trivial (\(\mu_{\mathrm{s}}=0\), equivalent to \(S\equiv 1\))";
  in this case there is a one-to-one correspondence between boundary density and point defect atomic data, which can be unambiguously recovered through the pole--residue structure of rational Herglotz.
\end{enumerate}
\end{proposition}

\begin{proof}[Proof outline]
The atomic part produces Lorentzian superposition as the discrete terms of the Poisson representation in Proposition \ref{prop:xi-phase-derivative-poisson-superposition};
for finite point defects this superposition gives rational Herglotz (hence real analytic boundary density) and can be uniquely recovered by pole--residue inversion (Theorem \ref{thm:xi-offcritical-poisson-herglotz-inversion}).
If \(\mu_{\mathrm{s}}\neq 0\), then its boundary contribution in the measure sense is a singular part, which cannot be identical to any absolutely continuous (let alone real analytic) density;
thus real analyticity forces \(\mu_{\mathrm{s}}=0\).
The reverse direction holds directly from the Lorentzian rational representation when \(S\equiv 1\) and finite point defects. This completes the proof.
\end{proof}

\begin{theorem}["Critical line concentration" and "no singular inner factor" combined criterion: RH-type sufficient condition template]\label{thm:xi-rh-sufficient-template-inner-clark-conservative}
In the above framework, if the following are satisfied:
\begin{enumerate}
  \item \(D\) is of bounded type on \(\mathbb{H}_{1/2}\) and satisfies reflection symmetry \eqref{eq:xi-inner-reflection-symmetry};
  \item Critical line pure-phase property \(|D(\tfrac12+it)|=1\) holds almost everywhere (thus \(D\) is an inner function, Theorem \ref{thm:xi-hardy-inner-characterization});
  \item For almost every \(\alpha\in\TT\), the Clark measure \(\sigma_\alpha\) is absolutely continuous (equivalent to trivial singular inner factor);
\end{enumerate}
then all non-trivial structure of \(D\) is determined by the Blaschke factor, and off-line zeros and phase velocity satisfy the invertible correspondence of Proposition \ref{prop:xi-phase-derivative-poisson-superposition}.
Furthermore, if one additionally imposes the right half-plane strict contraction regime induced by "commuting polar decomposition/joint unitary dilation realization of invisible generators" (e.g., Theorem \ref{thm:op-algebra-hzom-commuting-polar-forces-critical-line} or the contraction template of Proposition \ref{prop:xi-hzom-functional-equation}), then the no-zero event of \(D\) in \(\Re(s)>\tfrac12\), so that under reflection symmetry all its non-trivial zeros are forced to compress onto the critical line \(\Re(s)=\tfrac12\).
\end{theorem}

\endinput
